%% cap01-planteamiento.tex — Planteamiento del problema (versión sept-2026)

\chapter{Planteamiento del problema}
\label{cap:planteamiento}

\section{Contexto}
\label{sec:contexto}

En pocos años, la inteligencia artificial pasó a formar parte de la agenda educativa de
muchos gobiernos. Para 2021, más de treinta países habían publicado estrategias nacionales
de inteligencia artificial, y en ellas la educación aparece de manera casi constante
\citep[p.~527]{schiff2022education}. Desde entonces se han sumado planes de acción,
lineamientos y programas dedicados a la educación, como los que el gobierno chino y varias
de sus provincias publicaron entre 2025 y 2026.

Estos documentos no son todos del mismo tipo. La mayoría son estrategias generales de
inteligencia artificial en las que la educación ocupa algunas secciones, junto a la
industria, la salud o la seguridad. Otros son planes dedicados por completo a la educación.
En ambos casos, lo que un gobierno escribe sobre educación establece qué espera que se
enseñe sobre esta tecnología, cómo quiere que se use en las escuelas y qué pide a los
docentes.

Al leer varios de estos documentos, lo primero que se nota es que se parecen. Hablan de
competencias para el futuro, de ética, de formación docente, de innovación y de reducir
brechas. Ese parecido tiene explicaciones posibles: los gobiernos leen los documentos de
otros gobiernos, siguen las recomendaciones de los organismos internacionales y enfrentan
presiones parecidas. Lo que no dice ese parecido es si detrás de las mismas palabras hay
las mismas prioridades.

\section{Problema de investigación}
\label{sec:problema}

Usar el mismo lenguaje no significa proponer lo mismo. Los estudios que han comparado
políticas de inteligencia artificial lo muestran de dos maneras. En los principios éticos,
los documentos coinciden en unos cuantos términos, como transparencia, justicia o
privacidad, pero difieren en cómo los interpretan y en cómo piensan aplicarlos
\citep{jobin2019landscape, fjeld2020principled}. En educación, \citet{schiff2022education}
analizó 24 estrategias nacionales y encontró que todas hablaban de educación, pero casi
siempre para referirse a formar expertos y preparar a la fuerza laboral. Solo nueve
trataban con alguna profundidad el uso de la inteligencia artificial para enseñar y
aprender \citep[p.~535]{schiff2022education}.

Saber qué tan parecidas o distintas son estas propuestas exige compararlas tema por tema,
y las comparaciones disponibles tienen límites que impiden hacerlo. Comparan pocos casos,
entre dos y cuatro países \citep{kundu2025aipolicies, martinezrojas2026regulacion,
lara2026gobernanza}. Trabajan con documentos en inglés o traducidos al inglés, y sus
autores reconocen el sesgo que eso introduce \citep{jobin2019landscape,
schiff2022education}. Cada estudio usa categorías propias. Y la codificación la hace una
sola persona, de modo que no se puede saber cuánto dependen los resultados de quien leyó
\citep[p.~533]{schiff2022education}.

El problema de esta tesis es, entonces, que no sabemos qué tan parecido es lo que proponen
para la educación las políticas de inteligencia artificial de distintos países cuando se
comparan tema por tema, en sus idiomas originales y con un procedimiento que otra persona
pueda repetir.

\section{Preguntas de investigación}
\label{sec:preguntas}

\subsection{Pregunta central}

¿Qué tan parecido es lo que proponen para la educación las políticas de inteligencia
artificial de distintos países cuando se comparan tema por tema?

\subsection{Preguntas específicas}

\begin{enumerate}
  \item[\textbf{P1.}] ¿Qué temas educativos proponen las políticas de cada país, y plantean
  la educación como formación para la inteligencia artificial o como uso de la inteligencia
  artificial para enseñar y aprender?
  \item[\textbf{P2.}] Dentro de un mismo tema, ¿qué países proponen cosas cercanas y cuáles
  se alejan?
  \item[\textbf{P3.}] ¿Esas cercanías y distancias se mantienen cuando se controla el idioma
  de los documentos y el grado de acuerdo entre los modelos que los clasifican?
\end{enumerate}

La primera pregunta describe lo que propone cada política. La segunda compara esas
propuestas entre países. La tercera revisa si la comparación depende del idioma o del
instrumento con que se hizo.

\section{Objetivos}
\label{sec:objetivos}

\subsection{Objetivo general}

Comparar lo que proponen para la educación las políticas de inteligencia artificial de
distintos países, tema por tema, mediante procesamiento de lenguaje natural y un panel de
modelos de lenguaje multilingües.

\subsection{Objetivos específicos}

\begin{enumerate}
  \item[\textbf{OE1.}] Identificar los temas educativos que propone cada política y
  clasificarlos como formación para la inteligencia artificial o como uso de la
  inteligencia artificial para enseñar y aprender (P1).
  \item[\textbf{OE2.}] Medir, dentro de cada tema, la distancia entre lo que proponen los
  distintos países (P2).
  \item[\textbf{OE3.}] Contrastar esas distancias con un control de idioma y con el grado de
  acuerdo entre los modelos del panel (P3).
\end{enumerate}

\section{Justificación}
\label{sec:justificacion}

La justificación pedagógica está en el objeto de estudio. Lo que una política dice sobre
educación orienta decisiones concretas: qué contenidos de inteligencia artificial entran al
currículo, para qué niveles, qué se espera de los docentes y si la tecnología se piensa
como materia de estudio o como herramienta para enseñar. Comparar esas decisiones entre
países permite ver que no hay una sola manera de tomarlas.

La justificación metodológica está en el procedimiento. El procesamiento de lenguaje
natural permite comparar muchos documentos en sus idiomas originales, fragmento por
fragmento, algo que la lectura manual de una sola persona no puede hacer con consistencia.
Usar un panel de modelos en lugar de un solo codificador permite medir cuánto coinciden y,
con ello, cuánto confiar en cada resultado. Esta postura sigue a
\citet{grimmer2013text}: los métodos automáticos amplían la capacidad del analista, pero no
sustituyen la interpretación y deben validarse. El panel responde además a un riesgo
documentado: que la elección de un modelo o de una instrucción determine la conclusión
\citep{baumann2025llmhacking}.

La tesis también parte de mi trabajo. Diseño currículo de cursos sobre inteligencia
artificial en DeepLearning.AI, y estoy por concluir la edición 2026 del Programa regional
de formación en planeamiento y gestión de políticas educativas del IIPE UNESCO
\citep{iipe2026prf}. En el primer trabajo decido qué deben aprender las personas sobre esta
tecnología; en el segundo aprendo a analizar las políticas que orientan esas decisiones.
Esta tesis está en el punto donde ambos se cruzan.

\section{Alcances y limitaciones}
\label{sec:alcances}

El corpus reúne 16 documentos oficiales de ocho países: Alemania, Australia, Canadá, China,
Colombia, Estados Unidos, México y Sudáfrica. Están publicados entre 2020 y 2026 en alemán,
chino, español e inglés, y se analizan en su idioma original, divididos en 3{,}211
fragmentos. Siete de los países aportan un documento general de inteligencia artificial,
como estrategias nacionales, planes de acción o, en Sudáfrica, el informe de una comisión
presidencial; en el caso de México es el Plan Nacional de Inteligencia Artificial de abril de 2026
\citep{atdt2026plan}. China aporta nueve documentos, la mayoría dedicados a la educación.
Para que China no pese más que los demás en las comparaciones, se excluye de ellas su
15.º Plan Quinquenal, y esa exclusión se declara en cada resultado. Como el corpus mezcla
estrategias generales y planes educativos, lo que se compara son los fragmentos que tratan
de educación, y el tipo de documento del que proviene cada uno se registra y se reporta.

El estudio es exploratorio. Describe en qué se parecen y en qué se distinguen las
propuestas; no prueba por qué. Analiza lo que los documentos dicen, no si los gobiernos lo
cumplen: trabaja en el plano del texto y del discurso, no en el de los efectos
\citep{rizvi2010globalizing}. Con ocho casos no hay generalización estadística posible. En
Alemania, Australia, Canadá y Estados Unidos la educación depende en buena medida de los
estados o provincias, de modo que el documento nacional no equivale a la política que rige
en sus escuelas \citep{bray1995levels}.

Las categorías con las que se comparan los temas provienen de un marco internacional para
responsables de política (capítulo~\ref{cap:acercamiento}). Algunos países pudieron
redactar sus documentos siguiendo ese mismo marco. Por eso la tesis no mide cuánto se apega
cada país a él, sino qué tan lejos quedan los países entre sí dentro de cada tema: dos
países pueden tratar el mismo tema y proponer cosas muy distintas.
